\section{Theoretical Framework}
\label{sec:theory}

\subsection{Marchenko-Pastur Law for Honest Gradients}

\begin{definition}[Marchenko-Pastur Distribution]
Let $X \in \mathbb{R}^{n \times d}$ with i.i.d.\ zero-mean variance-$\sigma^2$ entries.
As $n,d \to \infty$ with $d/n \to \gamma \in (0,\infty)$, the empirical spectral distribution
of $\frac{1}{n}X^TX$ converges weakly to the \emph{Marchenko-Pastur law} with density:
\begin{equation}
\rho(\lambda) = \frac{1}{2\pi\sigma^2\gamma\lambda}
  \sqrt{(\lambda_+ - \lambda)(\lambda - \lambda_-)},
  \quad \lambda \in [\lambda_-,\, \lambda_+]
\end{equation}
where $\lambda_{\pm} = \sigma^2(1 \pm \sqrt{\gamma})^2$ are the bulk edges.
\end{definition}

Our core observation is that honest gradients, \emph{even under Non-IID heterogeneous data},
produce covariance matrices whose bulk eigenspectrum follows the MP law.
The key technical step is a decomposition lemma:

\begin{lemma}[Non-IID Decomposition]
\label{lem:noniid}
Under Assumption~2, each honest gradient decomposes as $g_i^t = \bar{g}^t + \Delta_i^t$,
where $\bar{g}^t = \frac{1}{|\mathcal{H}|}\sum_{i \in \mathcal{H}} g_i^t$ and
$\mathbb{E}[\Delta_i^t] = 0$, $\mathbb{E}[(\Delta_i^t)_j^2] \leq \sigma^2$.
The gradient matrix $G$ then has covariance
$C = \frac{1}{n}G^TG = \bar{g}(\bar{g})^T + \frac{1}{n}\sum_i \Delta_i^t(\Delta_i^t)^T + R$,
where $\|R\|_2 \leq 2\|\bar{g}\| \cdot \sigma/\sqrt{n}$ (cross term, vanishes as $n \to \infty$).
\end{lemma}

The rank-1 term $\bar{g}(\bar{g})^T$ contributes at most one outlier eigenvalue.
The remaining term $\frac{1}{n}\sum_i \Delta_i(\Delta_i)^T$ has zero-mean entries with
bounded second moment $\sigma^2$:

\begin{theorem}[MP Law for Honest Gradients]
\label{thm:mp_honest}
Under Assumption~2, as $n, d \to \infty$ with $d/n \to \gamma$, the \emph{bulk}
eigenvalue distribution of $C = \frac{1}{n}G^TG$ converges to the MP distribution
$\mathrm{MP}(\gamma, \sigma^2)$.
\end{theorem}

\begin{proof}[Proof sketch]
By Lemma~\ref{lem:noniid}, the bulk spectrum is governed by
$\frac{1}{n}\sum_i \Delta_i \Delta_i^T$.
Each $\Delta_i$ has i.i.d.-like coordinate structure (zero mean, bounded variance $\sigma^2$
coordinate-wise) but not necessarily Gaussian.
Tao and Vu's universality theorem~\cite{rmt_modern} establishes MP convergence under exactly
these conditions: bounded first four moments suffice for the limiting spectral distribution
to equal that of a Gaussian matrix with the same first two moments.
The rank-1 term $\bar{g}(\bar{g})^T$ creates a single spike eigenvalue above $\lambda_+$
when $\|\bar{g}\|^2 > \sigma^2 \sqrt{d/n}$ (BBP transition~\cite{bbp_transition}),
which is absorbed into the spectral filter as a detectable outlier but does not affect
the bulk convergence.
Hence the bulk of the empirical spectral distribution converges to MP$(\gamma, \sigma^2)$.
\end{proof}

\begin{remark}[Non-IID Concentration at Finite $n, d$]
The asymptotic MP convergence (large $n, d$) implies finite-sample concentration
at rate $O(1/\sqrt{\min(n,d)})$ by standard matrix concentration inequalities~\cite{rmt_modern}.
For our experimental setup ($n \geq 32$, $d \geq 22 \times 10^6$), the bulk is well-concentrated,
and the KS detection test calibrated against the fitted MP distribution is empirically validated
in the experiments (Section~\ref{sec:experiments}).
\end{remark}

\subsection{Byzantine Perturbations Create Spectral Anomalies}

\begin{theorem}[Spectral Anomaly Detection]
\label{thm:spectral_anomaly}
Let $f$ Byzantine clients inject arbitrary gradients, yielding perturbed matrix
$\tilde{G} = G + E$ where $E$ has $f$ non-zero rows.
If the Byzantine attack creates a rank-$f$ perturbation of magnitude
$\|\delta_{\mathrm{byz}}\| > \sigma\sqrt{d/n}$, then the eigenvalue spectrum of
$\tilde{C} = \frac{1}{n}\tilde{G}^T\tilde{G}$ exhibits $f$ outlier eigenvalues above
$\lambda_+$ with probability $\geq 1 - e^{-k/\log^2 k}$ for sketch size $k \geq \Omega(\log d)$.
\end{theorem}

\begin{proof}
The Byzantine perturbation yields
$\tilde{C} = C + \frac{1}{n}(G^TE + E^TG) + \frac{1}{n}E^TE$.
The rank-$f$ term $\frac{1}{n}E^TE$ contributes $f$ eigenvalues.
By the BBP phase transition~\cite{bbp_transition}: a rank-$r$ perturbation of magnitude
$\tau > \sigma^2\sqrt{n/d}$ produces $r$ eigenvalue spikes above $\lambda_+$.
The perturbation term $\frac{1}{n}G^TE$ adds $O(\sigma\|\delta_{\mathrm{byz}}\|/\sqrt{n})$
to the spectral norm, which is dominated by the rank-$f$ spike when
$\|\delta_{\mathrm{byz}}\| > \sigma\sqrt{d/n}$.
KS testing against the fitted MP distribution detects these outliers with the stated
probability by sub-Gaussian concentration of the empirical spectral distribution.
\end{proof}

\subsection{Phase Transition Characterization}

\begin{theorem}[Detection Phase Transition]
\label{thm:phase_transition_det}
Let $\beta = f/n$ be the Byzantine fraction.
For any Byzantine detection algorithm using only gradient information:
\emph{(i)} \emph{Below} $\sigma^2\beta^2 < 0.25$: detection succeeds with probability $> 1 - \delta$
(Theorem~\ref{thm:spectral_anomaly});
\emph{(ii)} \emph{Above} $\sigma^2\beta^2 \geq 0.25$: no algorithm can reliably distinguish
Byzantine from honest gradients (Theorem~\ref{thm:phase_impossibility}).
\end{theorem}

The threshold $\sigma^2\beta^2 = 0.25$ arises from the BBP
transition~\cite{bbp_transition}: a rank-$f$ perturbation with $\beta = f/n$ Byzantine fraction
exits the MP bulk if and only if its normalized magnitude exceeds $\sigma^2\sqrt{n/d}$,
which in the federated gradient setting yields the product condition $\sigma^2\beta^2 < 1/4$.
Equivalently, $\sigma\beta < 1/2$: the honest gradient noise amplitude must exceed twice
the Byzantine signal amplitude for detection to be information-theoretically feasible.
